\documentclass[../position-paper.tex]{subfiles}

\textbf{Philosophy of Computer Science}. Epistemological and methodological problems of computer science and of its objects of study: the computing machine and computation qua phenomenon. This includes the status of the Church-Turing-Post thesis as a thesis: what kind of claim it is, what evidence bears on it, and its role in the identity of the discipline.


\textbf{Metaphysics of Computing}. The ontological status of computational and virtual entities: whether computations, programs, algorithms, abstract machines, virtual objects, and simulated realities are real, and of what kind. Questions of implementation and the abstract/concrete relation, computational identity and individuation, the physical Church-Turing-Post thesis, and pancomputationalist claims about the physical world, as well as post-structuralist and Continental treatments of the virtual and the simulacrum.



\textbf{Theory, Logic, and Mathematics of Computation}. Axiomatic or categorical formulations of the notion of effective procedure or effective calculability, including the Church-Turing-Post thesis and the adequacy of a given axiomatization to the informal notion. Research in automated and interactive reasoning, as well as proof theory and proof mining.



\textbf{Sociopolitical and Aesthetic Problems of Computing}. Study of the political, social, cultural, and aesthetic dimensions of computing as a situated practice, with a strong historical focus. This includes the labour and material conditions of computing, its infrastructures and their political economy, the aesthetics of software and computational artifacts, and historiographical work that recovers underrepresented or marginalized traditions in the history of computation.



\textbf{Philosophy of Artificial Intelligence and Its Derivatives}. Critique and theorization of artificial intelligence systems, both as technical objects and as sociotechnical systems. This spans the conceptual foundations of machine learning and contemporary AI, questions of agency, understanding, and representation in such systems, their epistemic opacity, and the entanglement of technical architectures with the social, economic, and institutional arrangements in which they are built and deployed.


\textbf{Philosophy of Technology}. Philosophical and post-structuralist studies on the use of the machine metaphor in discourse in the age of accelerationism and techno-feudalism. This includes the genealogy and rhetoric of the machine as a figure of thought, critiques of technological determinism and solutionism, and analyses of how contemporary technological formations reshape labour, power, and subjectivity.



\textbf{Unconventional Computing}. Frontier theoretical research in super-Turing, biological, and quantum computing, including challenges to the physical Church–Turing-Post thesis posed by hypercomputation and physical models. This includes analog, relational, and membrane computing, the computational status of physical and living processes, and the conceptual question of what would count as computation beyond the classical model.


\textbf{Computation and Explanation}. Critical studies on the use of the computing-machine metaphor as explanation or justification in the sciences. This includes computationalism in the cognitive and life sciences, the explanatory role of computational and information-theoretic vocabulary, the difference between modelling a system computationally and claiming that it computes, and the limits of mechanistic explanation.


\textbf{General Systems Theory}. Research in the tradition of second-order cybernetics, enactive cognitive science, and non-mechanistic biologies that reject the computing machine as explanation. This includes autopoiesis and organizational closure, relational biology and the (M,R)-systems tradition, the epistemology of the observer in second-order cybernetics, and accounts of living and cognitive systems that resist reduction to effective computation